\documentclass[11pt]{article}
\usepackage{fontspec}
\usepackage[french]{babel}
\usepackage{amsmath}
\usepackage{amssymb}
\usepackage{amsthm}
\usepackage{mathrsfs}
\usepackage{tikz-cd}
\usepackage{enumitem}

\newcommand{\printedpage}[1]{\marginpar{\footnotesize p.~#1}}
\newcommand{\msp}[1]{\marginpar{\footnotesize [#1]}}
\newcommand{\Ch}{\underline{\mathrm{Ch}}}
\newcommand{\Ty}{\underline{\mathrm{Ty}}}
\newcommand{\Nerf}{\mathrm{Nerf}}

\begin{document}

\printedpage{39}
% head: A. Grothendieck : Dérivateurs, Ch. VII, Édition des Universités de Montpellier II et Paris VII
D'autre part, posant $l_i = \mathrm{long}(\pi_{i-1})$ (pour $i \in [1,r]$),
\[
f_i : \Delta_\varepsilon^{l_{i-1}} \longrightarrow \Delta_\varepsilon^{l_i\ (=\mathrm{long}(\pi_{i-1}))}
\]
est défini comme
\[
f_i = \mathfrak{f}_{\varepsilon,\pi_{i-1}(1)} * \mathfrak{f}_{\varepsilon,\pi_{i-1}(2)} * \cdots * \mathfrak{f}_{\varepsilon,\pi_{i-1}(l_{i-1})},
\]
de sorte que l'élément $(*)$ de $\Nerf_r(\Ty)$ est déterminé de façon unique par la donnée des $\varepsilon \in \{\pm 1\}$ et de
\[
\pi_\bullet = (\pi_i)_{i\in[0,r-1]} \in (\mathrm{Hom}_{\mathrm{Ens}}(\mathbf{N}^*,\mathbf{N}))^r
\]
satisfaisant les conditions $(**)$. Soit $F_{\varepsilon,\pi_\bullet}$ cet élément de $\Nerf_r(\Ty)$, et soit $N_r$ l'ensemble des $\pi_\bullet$ qu'on vient d'expliciter. Alors
\[
(F_{\varepsilon,\pi_\bullet})_{(\varepsilon,\pi_\bullet)\in\{\pm1\}\times N_r}
\]
est un système de générateurs libre du monoïde $\Nerf_r(\Ty)$.

Si $l_i = 0$ (i.e. \fbox{$\pi_{i-1} = 0$}), on doit avoir $\mathrm{supp}(\pi_{i+1}) = \emptyset$, i.e. $\pi_{i+1} = 0$, d'où $l_{i+1} = 0$, et de même pour les $l_j$ et $\pi_j$ suivants. Géométriquement l'hypothèse s'exprime par la condition que le but
\msp{page 58}
de $f_i$ est l'élément 1 de Ty (correspondant à l'intervalle ponctuel), et comme $\mathrm{Hom}(1,\tau) = \emptyset$ pour $\tau \neq 1$, cela signifie que les flèches qui suivent $f_i$ doivent être toutes égales à la flèche $\mathrm{id}_1 : 1 \longrightarrow 1$.

\textbf{Remarque.} Dans la description combinatoire de $\Ty$ et de $\Nerf_\bullet(\Ty)$, l'ensemble des `signatures' $\{\pm1\}$ joue un rôle assez bidon. Il peut être remplacé par n'importe quel ensemble $\Sigma$ (dit `ensemble des signatures'), et on construit alors une catégorie monoïdale (ou un monoïde catégorique) ayant comme ensemble d'objets le monoïde libre engendré par $\Sigma$, interprétés comme des intervalles d'entiers $[0,n]$, munis de `signatures' dans $\Sigma$ pour chaque intervalle de transition, ou ce qui revient au même, des `graphes' orientés $\Sigma$-pondérés de type combinatoire un intervalle $\boldsymbol{\Delta}_i$, ce qui revient à la donnée d'une suite (ordonnée) de $n$ éléments $\varepsilon_1,\ \varepsilon_2,\ \ldots\ ,\ \varepsilon_n \in \Sigma$, i.e. d'un `mot' $\varepsilon_1\varepsilon_2\ldots\varepsilon_n$, qu'on préfère écrire $\Delta_{\varepsilon_1}\Delta_{\varepsilon_2}\cdots\Delta_{\varepsilon_n}$. On peut définir de façon évidente les homomorphismes \textit{croissants} de tels graphes pondérés orientés (sous-entendu : avec l'élément source et l'élément but dans la structure, à respecter par les morphismes), sans avoir besoin d'une interprétation catégorique. La fonctorialité de l'opération de composition des intervalles pondérés, pour cette notion de morphisme, est évidente. Le
\msp{page 59}
choix $\Sigma = \{\pm1\}$ est à ma connaissance le seul utile, à cause de son lien avec la notion de chemins et le calcul de $\pi_0$ d'un objet de $X$ \texttt{[plutôt d'un objet} $X$ \texttt{de} Cat\texttt{]}, etc.
% foot: Version 22 décembre 2001

\printedpage{40}
% head: A. Grothendieck : Dérivateurs, Ch. VII, Édition des Universités de Montpellier II et Paris VII
\section*{7\quad Explicitation des catégories de fractions $X\Sigma^{-1}$ en termes des $\Ch(X;a,b)$}

Si $X$ est une catégorie, on désigne par $\Ch(X)$ la catégorie des chemins dans $X$ (cf. §4), et pour $a,b \in \mathrm{Ob}\,X$, par
\[
\Ch(X;a,b) \qquad \text{ou simplement } \Ch(a,b)
\]
celle des chemins de $a$ vers $b$, i.e. de source $a$, but $b$. L'intuition géométrique (qu'on explicitera), c'est qu'un tel chemin $c$ définit un chemin `direct' de $a$ vers $b$ (i.e. un chemin de type $\tau'$ `direct', i.e. toutes les flèches de transition sont directes, i.e. $\tau'$ de la forme $\Delta_+^n$), de même longueur que le chemin initial $c$, obtenu en inversant toutes les flèches rétrogrades dans le chemin $c$. Cette opération est effectivement possible dans $X$ lui-même si et seulement si lesdites flèches sont inversibles. On obtient alors une suite composable de flèches, dont le composé total (qu'on peut appeler `l'effet utile' du chemin)
\msp{page 60}
est une flèche de $a$ vers $b$. Dans le cas général, les chemins de $a$ vers $b$ peuvent être vus comme une sorte de `flèches généralisées' allant de $a$ vers $b$, dont la composition serait la composition évidente, et les identités les `chemins vides' relatifs aux $a \in \mathrm{Ob}\,X$. On trouve d'ailleurs bien ainsi une catégorie, ayant mêmes objets que $X$, mais où les
\[
\mathrm{Hom}_{\mathrm{Ch}}(a,b) = \mathrm{Ch}(a,b) = \mathrm{Ob}\,\Ch(a,b)
\]
sont prohibitivement gros. Cela s'exprime notamment par le fait que les seules flèches inversibles dans $X_{\mathrm{Ch}}$ sont les identités ! On n'a pas de foncteur naturel entre $X_{\mathrm{Ch}}$ et $X$ qui soit l'identité sur les objets. Car à un chemin $a \xrightarrow{c} b$ on ne peut associer de $X$-morphismes $a \longrightarrow b$ (sauf si $c$ \texttt{[est]} de type $\Delta_+$), et à un $X$-morphisme $a \xrightarrow{u} b$ on peut associer un chemin $c(u)$ de type $\Delta_+$, mais si $a \xrightarrow{u} b \xrightarrow{v} c$, $c(vu)$ est encore de type $\Delta_+$ alors que $c(v)c(u)$ est de type $\Delta_+^2$, ils ne sont pas égaux. Par contre, il y a une flèche naturelle entre ces deux éléments de $\Ch(a,c)$, flèche au dessus de l'unique homomorphisme
\[
\Delta_+ \xrightarrow{\mathfrak{f}_{+,2}} \Delta_+^2.
\]
\msp{page 61}
On a alors
\[
\mathfrak{f}_{+,2}^*(\underbrace{a \xrightarrow{u} b \xrightarrow{v} c}_{c(v)c(u)}) = (\underbrace{a \xrightarrow{vu} c}_{c(vu)}),
\]
d'où une flèche canonique
\[
c(vu) \longrightarrow c(v)c(u) \qquad \text{dans } \Ch(a,c).
\]
Cela montre que $c(vu)$ et $c(v)c(u)$ sont tout au moins dans une même composante connexe de $\Ch(a,c)$, ou (comme on dit encore) qu'ils sont `homotopes'. Cela suggère ici de modifier la définition de $\mathrm{Hom}_{\mathrm{Ch}}(a,b)$ en
\[
\pi(a,b) = \pi_0(\Ch(a,b)).
\]
% foot: Version 22 décembre 2001

\printedpage{41}
% head: A. Grothendieck : Dérivateurs, Ch. VII, Édition des Universités de Montpellier II et Paris VII
Le foncteur de composition des chemins
\[
\begin{array}{ccccc}
\Ch(a,b) & \times & \Ch(b,c) & \longrightarrow & \Ch(a,c)\\
(c & , & c') & \longmapsto & c' \circ c
\end{array}
\]
\texttt{[il y a trop de} $c$\texttt{]} (\footnote[17]{Ce foncteur est déduit des foncteurs
\[
\Ch_\tau(a,b) \times \Ch_{\tau'}(b,c) \longrightarrow \Ch_{\tau\tau'}(a,c)
\]
évidents, cf. page 77.}) induit une application sur les $\pi_0$, ce qui donne une loi de composition
\[
\pi(a,b) \times \pi(b,c) \longrightarrow \pi(a,c)
\]
sur les flèches. Il est immédiat qu'elle est associative (la composition des chemins l'étant), et unitaire. Elle admet comme flèche identique en $x$ soit le `lacet vide en $x$', lequel est (directement) homotope dans $\Ch(x,x)$ au chemin `dégénéré' $x \xrightarrow{\mathrm{id}} x$
\msp{page 62}
de type $\Delta_+$, ou aussi au chemin dégénéré de n'importe quel type $\tau$, disons
\[
x \xleftarrow{\mathrm{id}} x \xrightarrow{\mathrm{id}} x \xrightarrow{\mathrm{id}} x \xleftarrow{\mathrm{id}} x \xleftarrow{\mathrm{id}} x,
\]
un tel chemin n'étant autre que $p_\tau^*(x)$, où
\[
p_\tau : \tau \longrightarrow 1
\]
est le morphisme de dégénérescence (l'unique morphisme dans $\Ty$ vers 1, qui est objet final), d'où
\[
p_\tau^* : \Ch_1(x,x) = \{x\} \longrightarrow \Ch_\tau(x,x) \subseteq \Ch(x,x).
\]
La catégorie $\Pi(X)$ ainsi obtenu doit s'appeler le \textit{groupoïde fondamental} de $X$. En effet, % [sic] obtenu -> obtenue, agreement with « La catégorie »

\begin{enumerate}[label=\alph*)]
\item toutes les flèches de $\Pi(X)$ sont inversibles, i.e. $\Pi(X)$ est bien un groupoïde, et
\item le foncteur canonique
\[
X \longrightarrow \Pi(X)
\]
qui est l'identité sur les objets, et qui est défini sur les flèches par $u \longmapsto c(u)$ comme ci-dessus, est \textit{universel} (dans Cat) pour les foncteurs (i.e. les Cat-morphismes) de $X$ vers des \textit{groupoïdes.} (Plus généralement encore, il est universel pour les foncteurs $X \xrightarrow{F} Y$ tels que $F$ transforme toute flèche de $X$ en \texttt{[une]} flèche inversible.)
\end{enumerate}
% foot: Version 22 décembre 2001

\printedpage{42}
% head: A. Grothendieck : Dérivateurs, Ch. VII, Édition des Universités de Montpellier II et Paris VII
\msp{page 63}
Pour prouver que toute flèche dans $\Pi(X)$ est inversible, comme tout chemin non vide est composé de chemins `élémentaires' de type $\Delta_+$ ou $\Delta_-$, il suffit de prouver que ces chemins sont inversibles. En fait, ils se correspondent biunivoquement,
\[
\Ch_{\Delta_+}(a,b) \simeq \Ch_{\Delta_-}(b,a),
\]
en associant au chemin de $a$ vers $b$
\[
c(u) : a \xrightarrow{u} b \qquad \text{de type } \Delta_+
\]
le chemin
\[
c(u)^o : b \xleftarrow{u} a \qquad \text{de type } \Delta_-
\]
de $b$ vers $a$. Je dis que
\[
\begin{array}{rcl}
c(u)^o c(u) & \sim & 1_a\\
c(u) c(u)^o & \sim & 1_b
\end{array}
\]
\[
\left.
\begin{tikzcd}
 & a \arrow[r, "u"] & b & a \arrow[l, "u"'] \\
 {} & a \arrow[u, "\mathrm{id}"'] \arrow[r, "\mathrm{id}"] & a \arrow[u, "u"'] & a \arrow[u, "\mathrm{id}"'] \arrow[l, "\mathrm{id}"'] \\
 {} \arrow[u] & & a &
\end{tikzcd}
\right\}\
\begin{array}{l}
\text{on a un homomorphisme dans } \Ch_{\Delta_+\Delta_-}(a,a)\\
\text{de } 1_a \text{ dans } c(u)^o c(u).
\end{array}
\]
\[
\left.
\begin{tikzcd}
 & b & a \arrow[l, "u"'] \arrow[r, "u"] & b \\
 {} & a \arrow[u, "u"'] & a \arrow[u, "\mathrm{id}"'] \arrow[l, "\mathrm{id}"'] \arrow[r, "\mathrm{id}"] & a \arrow[u, "u"'] \\
 {} \arrow[u] & & a &
\end{tikzcd}
\right\}\
\begin{array}{l}
\text{itou dans } \Ch_{\Delta_-\Delta_+}(b,b),\\
\text{de } 1_b \longrightarrow c(u)c(u)^0 % [sic] c(u)^0 -> c(u)^o, digit 0 printed instead of the letter o used elsewhere
\end{array}
\]
\texttt{[plutôt} $1_b \longleftarrow c(u)c(u)^0$\texttt{, viz.} % [sic] c(u)^0 -> c(u)^o, digit 0 printed instead of the letter o used elsewhere
\[
\begin{tikzcd}
b \arrow[d, "\mathrm{id}"] & a \arrow[l, "u"'] \arrow[r, "u"] \arrow[d, "u"] & b \arrow[d, "\mathrm{id}"] \\
b & b \arrow[l, "\mathrm{id}"'] \arrow[r, "\mathrm{id}"] & b \ \texttt{].}
\end{tikzcd}
\]
% foot: Version 22 décembre 2001

\printedpage{43}
% head: A. Grothendieck : Dérivateurs, Ch. VII, Édition des Universités de Montpellier II et Paris VII
\msp{page 64}
Notons que, plus généralement, pour tout chemin
\[
c : a \longrightarrow b, \qquad \text{de type } \tau,
\]
on définit un \textit{chemin opposé}
\[
c^o : b \longrightarrow a \qquad \text{de type } \tau^o,
\]
où
\[
\tau \longmapsto \tau^o : \mathrm{Ty} \longrightarrow \mathrm{Ty}
\]
est l'antiïnvolution du monoide libre Ty qui échange les deux générateurs $\Delta_-$, $\Delta_+$, et qui est donc donné sur l'élément général par
\[
(\Delta_{\varepsilon_1}\Delta_{\varepsilon_2}\cdots\Delta_{\varepsilon_n})^o = \Delta_{-\varepsilon_n}\Delta_{-\varepsilon_{n-1}}\cdots\Delta_{-\varepsilon_1}.
\]
Par récurrence sur la longueur de $c$ (\footnote[18]{inutile, résulte directement de la formule $(*)$ \texttt{[de la]} page suivante, en décomposant $c$ en composés de chemins `élémentaires' de type $\Delta_+$ ou $\Delta_-$.}), en utilisant le résultat précédent, que $c(u)$, $c(u)^o$ sont inverses l'un de l'autre dans $\Pi(X)$, on trouve
\[
\left.
\begin{array}{rcl}
c^o c & = & \mathrm{id}_a\\
c c^o & = & \mathrm{id}_b
\end{array}
\right\}\text{dans } \Pi(X).
\]

\textbf{Exemple.} Si
\[
\begin{array}{rcl}
\tau & = & \Delta_+\Delta_-\Delta_+^2\Delta_-\\
\tau^o & = & \Delta_+\Delta_-^2\Delta_+\Delta_-,
\end{array}
\]
et si on a un chemin
\[
c : a_0 \xrightarrow{u_1} a_1 \xleftarrow{u_2} a_2 \xrightarrow{u_3} a_3 \xrightarrow{u_4} a_4 \xleftarrow{u_5} a_5
\]
de type $\tau$, on en déduit le chemin
\[
c^o : a_5 \xrightarrow{u_5} a_4 \xleftarrow{u_4} a_3 \xleftarrow{u_3} a_2 \xrightarrow{u_2} a_1 \xleftarrow{u_1} a_0
\]
de type $\tau^o$. L'opération $c \longmapsto c^o$ satisfait en
\msp{page 65}
effet \texttt{[pour]}
\[
\begin{array}{ccccc}
a & \xrightarrow{\gamma} & b & \xrightarrow{\gamma'} & c\\
c & \xrightarrow{\gamma'^o} & b & \xrightarrow{\gamma^o} & a
\end{array}
\]
la relation
\[
(\gamma'\gamma)^o = \gamma^o\gamma'^o, \tag*{$(*)$}
\]
% foot: Version 22 décembre 2001

\end{document}
